\section{Discussion}
\label{sec:discussion}

\para{Three compounding mechanisms.} Reading the paper-, reader-, and researcher-layer results jointly, the overall ``signaling effect of gate-keeping diminishes'' resolves into three formally distinct mechanisms with different rates. {\bf (1) Bayesian decay} (\secref{sec:results-paper}): under fixed capacity, mutual information and tail recall fall with $\Nsub$. This is the classical capacity-strain effect and is escapable by expanding capacity linearly with $\Nsub$. {\bf (2) Load-induced noise} (\secref{sec:results-empirical}): when capacity expands linearly but the reviewer pool does not, $\sigrev^2(\Nsub)$ grows; this re-injects the same $\contrast \to 0$ behavior through the back door. \iclr exhibits exactly this pattern. {\bf (3) Bean counting and reader budget} (\secref{sec:results-researcher}, \secref{sec:results-reader}): even if per-paper precision were preserved, the aggregator-level signal (publication count) and the reader-level allocator (which paper to read) both fail at finite $\Nsub$ because of fixed budgets---$\budget$ papers per reader, $T$ years per researcher.

\para{Why no policy is a free lunch.} The capacity-policy comparison in \figref{fig:capacity} reads, at first glance, as ``just expand capacity proportionally to volume.'' But the empirical \iclr calibration in \secref{sec:results-empirical} undercuts this: \iclr did expand capacity linearly, and the signal degraded anyway. The model attributes the degradation to load-induced noise. Conversely, holding capacity fixed preserves per-acceptance precision but produces vanishing recall on the true tail of papers and crushes researcher identifiability via mass tying at zero. {\bf The trilemma is real: per-paper precision, researcher identifiability, and reader-level relevance cannot all be preserved as $\Nsub$ grows without proportional growth in capacity \emph{and} in reviewer effort \emph{and} in reader attention.}

\para{What the empirical pattern says.} The $-0.32$ slope on $\log \Nsub$ for $\hcontrast / \stdrating$ is a quantitatively meaningful effect. From $2017$ to $2026$, the standardized signal at \iclr fell by $58\%$ even as the venue maintained its acceptance rate near the historical average. A reader who in $2017$ knew that a randomly chosen accepted paper was $2.55$ rating-SDs better than a randomly chosen rejected paper now sees only a $1.08$-SD gap. The implication for downstream uses of the venue label---attention allocation, hiring, citation---is that $2026$ \iclr ``acceptance'' carries less than half the information it did in $2017$.

\para{Connection to prior work.} The model recovers \citet{akerlof1970lemons}, \citet{spence1973job}, and \citet{long2023lemons} as the no-capacity-constraint limit (flat green line in \figref{fig:capacity}). The tail-recall metric reproduces \citet{heckman2019tyranny}'s observation that the top tier captures only about half of the most influential papers. The growth of $\sigrev$ we infer for \iclr 2026 sits in the same range as the Cortes--Lawrence $\neurips$ 2014 reproducibility result. The $-0.32$ log-linear slope provides quantitative content for the qualitative concerns voiced by \citet{mohanty2025crisis}.

\para{Surprises.} Two findings ran against our prior expectations. First, the posterior contrast $\contrast$ \emph{rises} in $\Nsub$ under fixed $\Kcap$, not falls. The naive form of the hypothesis would have predicted the opposite. Operationalizing the question as ``how much information does an acceptance label carry'' (mutual information / tail recall) rather than ``how different are accepted and rejected papers on average'' is essential for getting the right qualitative picture. Second, linear capacity rescues researcher identifiability \emph{completely} in our simulation ($\rho \approx 0.78$ flat in $M_r$), a stronger positive result than we expected.

\subsection{Limitations}
\label{sec:limitations}

\para{Modeling assumptions.} We assume reviewer scores are conditionally independent given $\quality$. Real reviews share author/institution/topic priors, which would inflate effective $\sigrev$ and accelerate the decay we report. Our $\quality$ is one-dimensional; real paper quality is multidimensional (novelty, rigor, impact), and our $\quality$ should be read as a projection onto a single useful axis. The researcher submission rate $S$ is fixed; in reality, $S(\ability)$ is increasing in ability~\citep{kwiek2026topperformers}, which would partly rescue identifiability under fixed $\Kcap$---softening but not reversing the H3 conclusion.

\para{Empirical reach.} The pre-$2017$ \iclr \openreview JSON files lack reviewer-rating fields, so we cannot measure dispersion before then. Eight years of usable rating data are enough to estimate a $-0.32 \pm 0.08$ slope but not for a fully saturated time series. Other AI venues are smaller and have shorter rating histories; they are consistent with the pattern but not pinned down by it. The trend results are observational---load-induced noise is the most plausible explanation for the empirical degradation, but coexisting contributors (paper-topic broadening, ML reviewer pool shifts, LLM-assisted reviewing~\citep{latona2024lottery}) cannot be ruled out.

\para{Reader trust collapse.} Our negative-experience feedback rule did not trigger collapse with $\tau_{\text{bad}} = -0.3$ because the average accepted paper has positive quality under both regimes. A more demanding threshold (\eg expecting top-quartile quality from a prestigious venue) would change the qualitative picture; that case is the natural next experiment.

\para{Generalization.} \iclr may not generalize to fields with different review structures (single blind, journal-with-revisions, post-publication review). Our model is venue-agnostic but our empirical anchor is one venue.

\subsection{Broader implications}

If the $-0.32$ log-linear decay is real and continues, two practical recommendations follow. {\bf For readers and committees}: the venue label should be treated as one signal among many, not a decisive gate---author identity, topic, and post-publication signals must do the rest of the work. {\bf For venues}: maintaining acceptance rate is not enough. Either the reviewer pool must grow proportionally to submissions (so $\sigrev$ stays constant), or the venue must accept that its certification function is decaying. Mechanism-design changes (rebuttal, AC discussion, transparent \openreview, structured rubrics) plausibly reduce $\sigrev^2(\Nsub)$ but their net effect is empirically open.
